\documentclass[11pt,a4paper]{article}

\usepackage[margin=16mm,top=14mm,bottom=16mm]{geometry}
\usepackage{fontspec}
\usepackage{xeCJK}
\usepackage{xcolor}
\usepackage{graphicx}
\usepackage{array}
\usepackage{booktabs}
\usepackage{tabularx}
\usepackage{enumitem}
\usepackage{titlesec}
\usepackage{fancyhdr}
\usepackage{lastpage}
\usepackage{needspace}
\usepackage{hyperref}
\usepackage[most]{tcolorbox}
\usepackage{listings}

\setmainfont{Segoe UI}
\setsansfont{Segoe UI Semibold}
\setmonofont{Courier New}[Scale=0.88]
\setCJKmainfont{Microsoft YaHei}
\setCJKsansfont{Microsoft YaHei}
\setCJKmonofont{Microsoft YaHei}

\definecolor{Ink}{HTML}{17252D}
\definecolor{Muted}{HTML}{5B6B73}
\definecolor{Teal}{HTML}{087E73}
\definecolor{TealSoft}{HTML}{E9F5F2}
\definecolor{Blue}{HTML}{265D8F}
\definecolor{BlueSoft}{HTML}{EDF3F8}
\definecolor{Amber}{HTML}{B56A00}
\definecolor{AmberSoft}{HTML}{FFF4DF}
\definecolor{Red}{HTML}{A33A3A}
\definecolor{RedSoft}{HTML}{FBECEC}
\definecolor{CodeBg}{HTML}{111A20}
\definecolor{CodeFg}{HTML}{E7F0F2}
\definecolor{Line}{HTML}{CAD5D9}

\hypersetup{colorlinks=true,linkcolor=Teal,urlcolor=Blue,pdfauthor={YOUTH Vision Team},pdftitle={NX163 视觉侦察程序指令手册}}
\setlength{\parindent}{0pt}
\setlength{\parskip}{4pt}
\setlist[itemize]{leftmargin=1.5em,itemsep=2pt,topsep=3pt}
\setlist[enumerate]{leftmargin=1.7em,itemsep=2pt,topsep=3pt}
\renewcommand{\arraystretch}{1.25}

\titleformat{\section}{\Large\bfseries\color{Ink}}{}{0pt}{}[\vspace{2pt}\color{Teal}\titlerule]
\titleformat{\subsection}{\large\bfseries\color{Ink}}{}{0pt}{}
\titlespacing*{\section}{0pt}{12pt}{7pt}
\titlespacing*{\subsection}{0pt}{9pt}{4pt}

\pagestyle{fancy}
\fancyhf{}
\fancyhead[L]{\small\color{Muted}NX163 \textbar{} 视觉侦察运行手册}
\fancyhead[R]{\small\color{Muted}2026-09-06}
\fancyfoot[L]{\small\color{Muted}USB SSH: 192.168.55.1}
\fancyfoot[R]{\small\color{Muted}\thepage/\pageref{LastPage}}
\renewcommand{\headrulewidth}{0.4pt}
\renewcommand{\footrulewidth}{0pt}
\setlength{\headheight}{13pt}

\lstdefinestyle{shell}{
  language=bash,
  basicstyle=\ttfamily\small\color{CodeFg},
  keywordstyle=\color{white}\bfseries,
  stringstyle=\color{cyan!65},
  commentstyle=\color{gray!55},
  showstringspaces=false,
  breaklines=true,
  breakatwhitespace=false,
  columns=fullflexible,
  keepspaces=true,
  upquote=true
}

\newtcblisting{commandbox}[1][]{
  enhanced,
  breakable,
  listing only,
  listing engine=listings,
  listing options={style=shell},
  colback=CodeBg,
  colframe=CodeBg,
  arc=2mm,
  boxrule=0pt,
  left=3mm,right=3mm,top=2mm,bottom=2mm,
  #1
}

\newtcolorbox{statusbox}[1][]{enhanced,breakable,colback=TealSoft,colframe=Teal,boxrule=0.7pt,arc=2mm,left=3mm,right=3mm,top=2mm,bottom=2mm,#1}
\newtcolorbox{notebox}[1][]{enhanced,breakable,colback=BlueSoft,colframe=Blue,boxrule=0.7pt,arc=2mm,left=3mm,right=3mm,top=2mm,bottom=2mm,#1}
\newtcolorbox{warnbox}[1][]{enhanced,breakable,colback=AmberSoft,colframe=Amber,boxrule=0.7pt,arc=2mm,left=3mm,right=3mm,top=2mm,bottom=2mm,#1}
\newtcolorbox{dangerbox}[1][]{enhanced,breakable,colback=RedSoft,colframe=Red,boxrule=0.8pt,arc=2mm,left=3mm,right=3mm,top=2mm,bottom=2mm,#1}

\newcommand{\tagpill}[2]{\tcbox[on line,boxrule=0pt,arc=1.5mm,colback=#1!12,coltext=#1,left=1.2mm,right=1.2mm,top=0.4mm,bottom=0.4mm]{\scriptsize\bfseries #2}}
\newcommand{\filepath}[1]{\texttt{\small\detokenize{#1}}}

\begin{document}

\begin{tcolorbox}[enhanced,colback=Ink,colframe=Ink,arc=2mm,boxrule=0pt,left=5mm,right=5mm,top=5mm,bottom=5mm]
  {\fontsize{23}{28}\selectfont\bfseries\color{white}NX163 视觉侦察程序指令手册}\par
  \vspace{2mm}
  {\large\color{white!78}比赛侦察、实时识别、结果面板、日志与受控融合启动}\par
  \vspace{3mm}
  \tagpill{cyan}{主机 nx163@nx163-desktop}\quad
  \tagpill{green}{当前 image 自启动}\quad
  \tagpill{orange}{无框录像 600 s}\quad
  \tagpill{blue}{更新 2026-09-06}
\end{tcolorbox}

\begin{statusbox}[title={\bfseries 当前实机状态}]
\begin{tabularx}{\linewidth}{>{\bfseries}p{33mm}X}
登录用户 & \texttt{nx163} \\
视觉目录 & \filepath{/home/nx163/youth-vision-runtime} \\
ROS2 目录 & \filepath{/home/nx163/uav_ros2_project} \\
开机服务 & \texttt{youth-vision.service}: \texttt{enabled + active} \\
默认模式 & 比赛图案侦察 \texttt{competition\_selected.sh image} \\
录像 & 原始无框视频，1440x1080，60 FPS，最长 600 秒 \\
\end{tabularx}
\end{statusbox}

\begin{warnbox}[title={\bfseries 先看这三条}]
\begin{enumerate}
  \item 同一时间只运行一套占用相机的视觉程序。
  \item 结果面板只读文件，可以和比赛侦察同时运行。
  \item \texttt{competition\_selected.sh} 不启动飞控；\texttt{start\_single\_target\_fusion.sh} 会启动真实飞控融合链，二者不能混用。
\end{enumerate}
\end{warnbox}

\section{连接与身份确认}

\begin{commandbox}
ssh nx163@192.168.55.1
whoami
hostname
pwd
\end{commandbox}

正常应分别看到 \texttt{nx163}、\texttt{nx163-desktop}、\filepath{/home/nx163}。

\section{比赛侦察}

这是正式比赛选择逻辑的入口。它启动 Pose、分类、坐标解算、一套 MAVROS 和比赛选择器，并保存所有侦察结果。图案模式选择三个非空目标中价值最高者；数字模式选择三个非空数字的中位数目标。最终只生成一个候选坐标接口，但不会自行启动飞控。

\subsection{图案模式}
\begin{commandbox}
/home/nx163/youth-vision-runtime/scripts/competition_selected.sh image
\end{commandbox}

\subsection{数字模式}
\begin{commandbox}
/home/nx163/youth-vision-runtime/scripts/competition_selected.sh digit
\end{commandbox}

若提示已有侦察进程，先执行：
\begin{commandbox}
sudo systemctl stop youth-vision
/home/nx163/youth-vision-runtime/scripts/stop_recon_pipeline.sh
\end{commandbox}

\begin{notebox}[title={\bfseries 比赛输出}]
全部结果位于 \filepath{/home/nx163/youth-vision-runtime/recon_results/latest/target_*/result.json}；唯一选择写入 \filepath{/home/nx163/youth-vision-runtime/recon_results/latest/competition_selected.json}，并发布到 \texttt{/vision/competition\_selected\_target}。
\end{notebox}

\section{普通侦察与识别}

第二个参数是期望标靶数，范围为 0--20。\texttt{0} 表示持续运行，不因目标数量自动结束。

\textbf{单标靶数字侦察}
\begin{commandbox}
/home/nx163/youth-vision-runtime/scripts/start_recon_pipeline.sh digit 1
\end{commandbox}

\textbf{单标靶图案侦察}
\begin{commandbox}
/home/nx163/youth-vision-runtime/scripts/start_recon_pipeline.sh image 1
\end{commandbox}

\textbf{持续图案侦察}
\begin{commandbox}
/home/nx163/youth-vision-runtime/scripts/start_recon_pipeline.sh image 0
\end{commandbox}

普通侦察包含 Pose、分类、MAVROS 与坐标节点，但不启动比赛“三选一”选择器，也不启动飞控。

\subsection{停止侦察}
\begin{commandbox}
/home/nx163/youth-vision-runtime/scripts/stop_recon_pipeline.sh
\end{commandbox}

\begin{warnbox}
不要给停止脚本添加 \texttt{--help}。该脚本不解析帮助参数，执行后会直接停止当前侦察。
\end{warnbox}

\Needspace{72mm}
\section{网页程序}

\subsection{侦察结果面板}

结果面板不占用相机，可与比赛侦察同时运行。这里显式使用 6 米展示去重半径：

\begin{commandbox}
RADIUS_M=6.0 /home/nx163/youth-vision-runtime/scripts/start_recon_results_dashboard.sh
\end{commandbox}

\begin{center}
\href{http://192.168.55.1:8092/}{\texttt{http://192.168.55.1:8092/}}
\end{center}

\begin{commandbox}
/home/nx163/youth-vision-runtime/scripts/stop_recon_results_dashboard.sh
\end{commandbox}

\subsection{实时摄像头识别网页}

此程序独占相机。先停止开机侦察，再启动需要的模式：

\begin{commandbox}
sudo systemctl stop youth-vision
/home/nx163/youth-vision-runtime/scripts/stop_recon_pipeline.sh

WEB_PORT=8000 /home/nx163/youth-vision-runtime/scripts/start_youth_pipeline.sh image
\end{commandbox}

数字模式使用：
\begin{commandbox}
WEB_PORT=8000 /home/nx163/youth-vision-runtime/scripts/start_youth_pipeline.sh digit
\end{commandbox}

访问 \href{http://192.168.55.1:8000/}{\texttt{http://192.168.55.1:8000/}}。停止实时网页识别：

\begin{commandbox}
/home/nx163/youth-vision-runtime/scripts/stop_youth_pipeline.sh
\end{commandbox}

\subsection{普通侦察加实时网页}

以下入口同时启动普通侦察和实时网页，但不包含比赛选择器：

\begin{commandbox}
WEB_PORT=8000 /home/nx163/youth-vision-runtime/scripts/start_recon_web.sh image
\end{commandbox}

数字模式使用：
\begin{commandbox}
WEB_PORT=8000 /home/nx163/youth-vision-runtime/scripts/start_recon_web.sh digit
\end{commandbox}

访问 \href{http://192.168.55.1:8000/index_recon.html}{\texttt{http://192.168.55.1:8000/index\_recon.html}}。

\Needspace{68mm}
\section{开机服务与恢复}

当前开机服务固定启动比赛 \texttt{image} 模式。

\textbf{状态、启动与停止}
\begin{commandbox}
systemctl status youth-vision --no-pager
sudo systemctl start youth-vision
sudo systemctl stop youth-vision
\end{commandbox}

\textbf{开机自启动控制}
\begin{commandbox}
sudo systemctl enable youth-vision
sudo systemctl disable youth-vision
\end{commandbox}

\Needspace{25mm}
\textbf{持续查看本次开机日志}
\begin{commandbox}
journalctl -u youth-vision -b -f
\end{commandbox}

推荐使用下面的完整恢复顺序，避免旧子进程占用相机或 MAVROS：

\begin{commandbox}
sudo systemctl stop youth-vision
/home/nx163/youth-vision-runtime/scripts/stop_recon_pipeline.sh
sudo systemctl start youth-vision
\end{commandbox}

\section{结果、日志与录像}

\begin{tabularx}{\linewidth}{>{\bfseries}p{34mm}X}
当前侦察结果 & \filepath{/home/nx163/youth-vision-runtime/recon_results/latest/} \\
侦察日志 & \filepath{/home/nx163/youth-vision-runtime/logs/recon/} \\
识别事件 & \filepath{/home/nx163/youth-vision-runtime/logs/recognition/} \\
无框原始录像 & \filepath{/home/nx163/camera_recordings/} \\
服务定义 & \filepath{/etc/systemd/system/youth-vision.service} \\
\end{tabularx}

\textbf{当前会话目录}
\begin{commandbox}
readlink -f /home/nx163/youth-vision-runtime/recon_results/latest
\end{commandbox}

\textbf{最近录像与视觉日志}
\begin{commandbox}
ls -lht /home/nx163/camera_recordings | head
ls -lt /home/nx163/youth-vision-runtime/logs/recon | head
\end{commandbox}

\textbf{检查关键进程}
\begin{commandbox}
pgrep -af 'youth_vision_runner|mavros_node|recon_geolocator_node|competition_selector.py'
\end{commandbox}

\Needspace{92mm}
\section{飞控相关入口}

\begin{dangerbox}[title={\bfseries 受控测试区命令}]
以下程序会进入飞控、任务上传、模式切换或载荷相关链路。路径已经迁移为 NX163，但本次只完成静态依赖检查，没有执行真实飞行闭环。必须在拆桨、无载荷或经批准的 SITL/DRY\_RUN 环境完成首轮验证。
\end{dangerbox}

\subsection{读取配置并启动飞控 bringup}
\begin{commandbox}
cat /home/nx163/uav_ros2_project/config/recognition_mode.conf
bash /home/nx163/uav_ros2_project/scripts/start_uav.sh
\end{commandbox}

配置中的识别模式应为 \texttt{digit}、\texttt{image} 或 \texttt{none}。

\subsection{单标靶视觉飞控融合}
\begin{commandbox}
/home/nx163/uav_ros2_project/scripts/start_single_target_fusion.sh digit 90
/home/nx163/uav_ros2_project/scripts/start_single_target_fusion.sh image 90
\end{commandbox}

第二个参数为 \texttt{heading\_deg}。该入口会启动飞控 bringup、侦察、桥接和任务管理，不是纯识别命令。

\subsection{人工坐标输入}
\begin{commandbox}
bash /home/nx163/uav_ros2_project/scripts/start_manual_target.sh LATITUDE LONGITUDE HEADING_DEG
\end{commandbox}

\section{快速排查}

\begin{tabularx}{\linewidth}{>{\bfseries\color{Ink}}p{38mm}X}
提示 already running & 先停止 \texttt{youth-vision}，再运行 \texttt{stop\_recon\_pipeline.sh}。 \\
相机打不开 & 检查是否存在第二个 \texttt{youth\_vision\_runner} 或实时网页识别进程。 \\
UDP 15001 被占用 & 使用 \texttt{ss -lunp | grep 15001} 检查重复 MAVROS。 \\
有识别但无坐标 & 检查 FCU telemetry、RTK \texttt{fix\_type >= 6}、姿态与相对高度。 \\
录像没有生成 & 检查 \filepath{/home/nx163/camera_recordings/} 剩余空间是否低于 2 GiB。 \\
TensorRT 跨设备警告 & 与用户名路径无关；应在 NX163 本机重新生成或核验三个 engine。 \\
\end{tabularx}

\begin{notebox}[title={\bfseries 本版核验边界}]
已核验 NX163 路径、脚本语法、ROS2 包、比赛 image 自启动、单实例进程、600 秒参数链路和无框录像创建。未执行真实飞控、航点上传、模式切换及载荷动作。
\end{notebox}

\end{document}
